\chapter{Introduction}
\label{ch:intro}

\section{Motivation}

This report is a master's-thesis-style account of a staged modelling campaign for \emph{free-radical photopolymerization of poly(ethylene glycol) diacrylate} (PEGDA).
The engineering context is continuous-flow digital lithography of multi-porosity hydrogel patches, as envisaged in the DFG proposal \emph{Continuous-Flow Digital Lithography of Multi-Porosity Hydrogel Patches} (AVT.CVT, RWTH Aachen).
The scientific problem is older and more elementary than that device: given a known light history, and later a known flow, can one predict how many acrylate groups have reacted, whether a sample-spanning network exists, and how stiff the gel is---without treating those three questions as the same number?

A machine operator can print without answering them.
A working curve relating cure depth to the logarithm of dose \citep{Jacobs1992book} is a calibration for a given resin, lamp, and speed.
It is not a constitutive law.
Two parcels of resin can receive the same
\begin{equation}
  E = \int I(t)\,\mathrm{d}t
\end{equation}
and still differ in conversion $p$, gelation, and modulus, because bimolecular termination and oxygen inhibition make the outcome depend on the \emph{shape} of $I(t)$ \citep{Wydra2014,Lin2019,Dobson2024}.
A flowing jet continuously replenishes oxygenated liquid \citep{Slutzky2019}.
A DMD pixel is not a top-hat of light \citep{Montgomery2022,Meenakshisundaram2020}.
A dilute PEGDA solution may never gel even at full acrylate conversion, because too many additions close loops \citep{Zhu2020}.

The campaign therefore refuses a single mixed three-dimensional code on day one.
It implements seven stages (0--6), each with a frozen configuration, a forward run, named pass/fail gates, and one evidence tag (\texttt{proposed}, \texttt{smoke-tested}, \texttt{validated}, \texttt{matched}, or \texttt{invalid}).
As of 13 September 2026, Stages~0--6 are \tagvalid{} internally.
None is \tagmatch{} to FTIR, a camera point-spread function, an experimental critical speed, or rheology.

\section{Locked scientific object and non-goals}

\paragraph{Object.}
PEGDA \emph{chain-growth} free-radical photopolymerization: prescribed irradiance $I(t)$ maps to conversion $p$; later a DMD surface field and Beer--Lambert attenuation; later plug flow; then a loop-aware network post-processor $\eta(p,[\mathrm{PEGDA}]_0)$ and phantom modulus $G$.

\paragraph{Non-goals (parked).}
The full multi-species free-volume model of \citet{Dobson2024} as an everyday right-hand side; Navier--Stokes coupled to gelation; scattering optics; inverse DMD design; PNIPAM and thiol--ene closures; mixing eosin~Y Type~II bleaching ODEs into an Irgacure Type~I rate law; using a Jacobs working curve as the polymer state.

\paragraph{Two velocities.}
Pattern velocity $\mathbf{v}_{\mathrm{pattern}}$ (scan or mask) is independent of resin velocity $\mathbf{u}$.
A fluid element sees $\mathbf{v}_{\mathrm{pattern}}-\mathbf{u}$ \citep{Meenakshisundaram2020,Slutzky2019}.

\section{Intended reader}

The text is written for a starting graduate student who knows physical chemistry and transport but need not have worked on vat photopolymerization.
Jargon is defined once.
Evidence tags are used in the same sense throughout: engineering success (the solver finished) is not scientific validity.

\section{How this report is organised}

Chapter~\ref{ch:background} is the ontology: species, rates, light, oxygen, flow, and network, and the list of quantities that must not be identified.
Chapter~\ref{ch:literature} is a literature study of the five campaign papers and the supporting oxygen, reciprocity, and gelation literature, including the gap that the code was written to close.
Chapter~\ref{ch:methods} describes the Python package, the campaign quartet, numerical methods, and parameter books.
Chapter~\ref{ch:results} reports Stages~0--6 with gates, figures, and what each stage is \emph{not}.
Chapter~\ref{ch:matching} states what would be required to retag any stage \tagmatch.
Chapter~\ref{ch:conclusions} summarises findings, limitations, and parked work.

\section{Statement of originality and limitations}

No new laboratory experiment is reported.
The contribution is a citable modelling stack with separate chemistry keys, internal gates, and an explicit refusal to mix recipes.
Where a paper figure is reproduced only qualitatively (Montgomery Fig.~3-style Gaussians; Slutzky Figs.~2--3 illustrations), the tag remains \tagvalid, not \tagmatch.
Publisher PDFs in \texttt{Literature/} are local-only and are not redistributed.
